% !PRINTED 347
\weberpage{347}
%<<CONT
sche Näherungsmethode muss auf sie führen. Man bekommt nun die Gleichungen, deren Wurzeln die zweite, vierte, achte \ldots\ Potenz der Gleichung (4) sind:
\[
\begin{aligned}
x^3 - 4x^2 + 4x - 4 &= 0 \\
x^3 - 8x^2 - 16x - 16 &= 0 \\
x^3 - 96x^2 - 16^2 &= 0
\end{aligned}
\]
und
\[
x = \sqrt[8]{96} = 1{,}7692\ldots
\]
ist bereits ein in den vier ersten Decimalen genauer Werth. Der nächste Schritt würde kein anderes Resultat ergeben (weil die erste Potenz der Unbekannten in der letzten Gleichung fehlt); aber der darauf folgende ergiebt den in der sechsten Decimale genauen Werth
\[
x = \sqrt[32]{85032960} = 1{,}769293\ldots
\]

Wir wollen noch ein zweites Beispiel einer Gleichung fünften Grades betrachten, das Gelegenheit bietet, mehrere der Sätze des vorigen Paragraphen anzuwenden. Das Beispiel ist, wie die früheren, der Theorie der complexen Multiplication der elliptischen Functionen entnommen. Es ist die Gleichung fünften Grades
\weq{5}{x^5 - x^3 - 2x^2 - 2x - 1 = 0.}
Wir leiten daraus die Gleichung ab, deren Wurzeln die Quadrate der Wurzeln von (5) sind, indem wir in
\[
(x^5 - x^3 - 2x)^2 - (2x^2 - 1)^2 = 0
\]
$x^2$ durch $x$ ersetzen. Dies giebt
\weq{6}{x^5 - 2x^4 - 3x^3 - 1 = 0,}
und wenn wir dasselbe Verfahren zunächst noch zweimal anwenden, so erhalten wir die Gleichungen, deren Wurzeln die vierten und achten Potenzen der Wurzeln von (5) sind:
\weq{7}{x^5 - 10x^4 + 9x^3 - 4x^2 - 1 = 0,}
\weq{8}{x^5 - 82x^4 + x^3 - 36x^2 - 8x - 1 = 0.}
Die letzte dieser Gleichungen eignet sich zur Discussion. Sie hat nach dem \gesperrt{Cartesischen} Lehrsatz (\S.\,101) wenigstens eine positive Wurzel.

Wenn wir für die Gleichung (8) die in dem \gesperrt{Newton'schen} Kriterium (\S.\,101) vorkommenden Functionen
%CONT>>
